\section{Graphiques statistiques}

\begin{cadretheorie}
Un graphique statistique permet de rendre visibles les informations contenues dans une série de données.

Le choix du graphique dépend du type de variable étudiée.
\end{cadretheorie}

\begin{cadregris}{Choisir un graphique}
\begin{tabularx}{\textwidth}{>{\bfseries}p{4cm}X}
\toprule
Situation & Graphique adapté \\
\midrule
Caractère qualitatif & Diagramme en bâtons ou en secteurs \\
Caractère quantitatif discret & Diagramme en bâtons \\
Caractère quantitatif continu & Histogramme \\
Comparaison de distributions & Boîte à moustaches \\
Fréquences cumulées & Polygone des fréquences cumulées \\
\bottomrule
\end{tabularx}
\end{cadregris}